\chapter{几何多重网格、代数多重网格与源码实现}
\label{ch:software}

\begin{keybox}
\textbf{本章不把 AMReX、hypre、PETSc 等软件当作 API 名词表，而是直接从源码恢复它们的求解器结构。}
阅读重点固定为四个问题：
\begin{equation}
\boxed{
\text{谁构造 hierarchy}
\rightarrow
\text{谁保存 level state}
\rightarrow
\text{谁驱动 cycle}
\rightarrow
\text{每个数学操作最终调用什么 kernel/object}.
}
\label{eq:software-source-reading-chain}
\end{equation}
AMReX MLMG 用来观察 block-structured GMG；hypre BoomerAMG 用来观察分布式 AMG；PETSc PCMG/PCGAMG 用来观察对象组合式求解器架构。读完本章以后，读者应能从一个成熟求解器的源码入口一路追到 smoother、restriction、coarse operator、bottom solve 与 coarse-process reduction，而不是只会配置参数。
\end{keybox}

\section{源码快照与阅读边界}

为了避免源码随版本变化导致函数位置漂移，本章按固定源码快照讨论：

\begin{center}
\small
\begin{tabularx}{0.98\textwidth}{>{\bfseries}p{2.0cm}p{3.1cm}X}
\toprule
软件 & 源码快照 & 本章重点文件\\
\midrule
AMReX &
\texttt{fca855037a2184aae8fbfe4ff7b5435e182e7cfa} &
\texttt{Src/LinearSolvers/MLMG/AMReX\_MLMG.H}，
\texttt{AMReX\_MLLinOp.H}\\
hypre &
\texttt{5b582610a113ede9162a5744d33873b860498a11} &
\texttt{src/parcsr\_ls/par\_amg\_setup.c}，
\texttt{par\_cycle.c}，
\texttt{par\_strength.c}，
\texttt{par\_rap.c}\\
PETSc &
\texttt{0239ddb405da990d3e3adce5fef62d8582258e84} &
\texttt{include/petsc/private/pcmgimpl.h}，
\texttt{src/ksp/pc/impls/mg/mg.c}，
\texttt{gamg/gamg.c}\\
\bottomrule
\end{tabularx}
\end{center}

对应源码可从以下固定链接进入：

\begin{itemize}
\item \href{https://github.com/AMReX-Codes/amrex/blob/fca855037a2184aae8fbfe4ff7b5435e182e7cfa/Src/LinearSolvers/MLMG/AMReX_MLMG.H}{AMReX \texttt{AMReX\_MLMG.H}}
\item \href{https://github.com/AMReX-Codes/amrex/blob/fca855037a2184aae8fbfe4ff7b5435e182e7cfa/Src/LinearSolvers/MLMG/AMReX_MLLinOp.H}{AMReX \texttt{AMReX\_MLLinOp.H}}
\item \href{https://github.com/hypre-space/hypre/blob/5b582610a113ede9162a5744d33873b860498a11/src/parcsr_ls/par_amg_setup.c}{hypre \texttt{par\_amg\_setup.c}}
\item \href{https://github.com/hypre-space/hypre/blob/5b582610a113ede9162a5744d33873b860498a11/src/parcsr_ls/par_cycle.c}{hypre \texttt{par\_cycle.c}}
\item \href{https://github.com/petsc/petsc/blob/0239ddb405da990d3e3adce5fef62d8582258e84/include/petsc/private/pcmgimpl.h}{PETSc \texttt{pcmgimpl.h}}
\item \href{https://github.com/petsc/petsc/blob/0239ddb405da990d3e3adce5fef62d8582258e84/src/ksp/pc/impls/mg/mg.c}{PETSc \texttt{mg.c}}
\item \href{https://github.com/petsc/petsc/blob/0239ddb405da990d3e3adce5fef62d8582258e84/src/ksp/pc/impls/gamg/gamg.c}{PETSc \texttt{gamg.c}}
\end{itemize}

本章不是逐行注释整个项目。我们只追能够解释求解器骨架的调用链，并把它重新映射到前面已经建立的
\[
\text{state}
\rightarrow
\text{DAG}
\rightarrow
\text{ownership}
\rightarrow
\text{machine cost}
\]
心智模型。

\section{GMG 与 AMG 的实现边界}

几何多重网格与代数多重网格的根本区别发生在 \textbf{setup}。

GMG 已知几何层级，因此可以直接构造
\begin{equation}
\Omega_h
\rightarrow
\Omega_H,
\qquad
R,\ P,\ A_H.
\end{equation}

AMG 的输入通常是稀疏矩阵 $A$。它必须先从矩阵恢复粗空间：
\begin{equation}
\boxed{
A
\rightarrow
S\ \text{(strength graph)}
\rightarrow
\text{coarsening/aggregation}
\rightarrow
P,R
\rightarrow
A_c=RAP.
}
\label{eq:software-amg-setup-chain}
\end{equation}

这一差别在源码里非常明显。AMReX 的主要复杂度集中在几何层级、BoxArray、DistributionMapping 与离散算子；BoomerAMG 的 setup 则直接出现 strength matrix、C/F splitting、interpolation builder 和 ParCSR RAP。

AMG 的两个基础复杂度指标仍然值得保留。网格复杂度定义为
\begin{equation}
C_g
=
\frac{\sum_{\ell=0}^{L}n_\ell}{n_0},
\label{eq:grid-complexity}
\end{equation}
算子复杂度定义为
\begin{equation}
C_o
=
\frac{\sum_{\ell=0}^{L}\operatorname{nnz}(A_\ell)}
{\operatorname{nnz}(A_0)}.
\label{eq:operator-complexity}
\end{equation}
前者衡量整个 hierarchy 的未知量规模，后者进一步计入粗层矩阵的稠密化。后面会看到，PETSc GAMG 甚至直接在 setup 源码中统计并打印 operator complexity。

\section{AMReX MLMG 的求解器结构}
\label{sec:software-amrex-source}

AMReX MLMG 最值得学习的不是某个 stencil，而是它把\textbf{求解器控制流}和\textbf{离散算子实现}分成两个对象：

\begin{equation}
\boxed{
\texttt{MLMGT}
=
\text{cycle / convergence / temporary-state driver}
}
\end{equation}
\begin{equation}
\boxed{
\texttt{MLLinOpT}
=
\text{operator / smooth / restriction / interpolation / hierarchy metadata}.
}
\end{equation}

这是理解整个源码的第一把钥匙。

\subsection{\texttt{solve()} 的外层求解控制}

从 \texttt{MLMGT<MF>::solve()} 进入，可以看到核心顺序近似为

\begin{verbatim}
prepareForSolve(a_sol, a_rhs);
computeMLResidual(finest_amr_lev);
...
for (int iter = 0; iter < niters; ++iter) {
    oneIter(iter);
    computeResidual(finest_amr_lev);
    ...
}
\end{verbatim}

源码并没有在 \texttt{solve()} 里展开 stencil。它做的是：

\begin{enumerate}
\item 根据输入 solution/RHS 建立 solve state；
\item 计算初始 composite residual；
\item 做全局 residual/RHS norm reduction；
\item 反复调用 \texttt{oneIter()}；
\item 每轮以后重新计算 residual 并判断收敛。
\end{enumerate}

因此 \texttt{solve()} 对应的是\chapref{ch:parallel-model}中的\textbf{最外层控制 DAG}，不是局部数值 kernel。

源码还有一个很值得注意的 GPU 细节：当调用
\texttt{setNoGpuSync(true)} 时，\texttt{solve()} 用
\texttt{Gpu::SyncAtExitOnly} 与
\texttt{Gpu::SingleStreamRegion}
包住整段求解。这说明 AMReX 已经把“函数调用结束是否隐式同步”视为求解器级策略，而不是每个 kernel 自己决定。

\subsection{\texttt{oneIter()} 的 AMR/MG 层级控制}

AMReX MLMG 有两个不同的层级索引：

\begin{equation}
\boxed{
\texttt{alev}
=
\text{AMR level},
\qquad
\texttt{mglev}
=
\text{某个 AMR level 内部的 multigrid level}.
}
\label{eq:software-amrex-two-hierarchies}
\end{equation}

\texttt{oneIter()} 先沿 AMR level 从 fine 向 coarse 推进：
\begin{verbatim}
for (int alev = finest_amr_lev; alev > 0; --alev) {
    miniCycle(alev);
    LocalAdd(sol[alev], cor[alev][0], ...);
    computeResWithCrseSolFineCor(alev-1, alev);
    ...
}
\end{verbatim}

而 \texttt{miniCycle(alev)} 的主体只是
\begin{verbatim}
mgVcycle(amrlev, 0);
\end{verbatim}

也就是说：

\begin{equation}
\boxed{
\texttt{oneIter}
\text{ 管 AMR 复合层级，}
\qquad
\texttt{mgVcycle}
\text{ 管单个 AMR level 内部的 MG 层级。}
}
\end{equation}

这是阅读 AMReX MLMG 时最容易混淆的结构。若把 \texttt{alev} 和 \texttt{mglev} 都理解成“粗网格层”，后面的 coarse-fine residual 与 correction 逻辑就会完全看乱。

\subsection{\texttt{mgVcycle()} 的 V-Cycle 实现}

AMReX 的 \texttt{mgVcycle()} 很适合源码精读，因为它与标准 V-cycle 一一对应。

下降阶段的核心源码是：

\begin{verbatim}
setVal(cor[amrlev][mglev], 0.0);
linop.smooth(..., cor[amrlev][mglev],
             res[amrlev][mglev], ..., nu1);

computeResOfCorrection(amrlev, mglev);

linop.restriction(amrlev, mglev+1,
                  res[amrlev][mglev+1],
                  rescor[amrlev][mglev]);
\end{verbatim}

数学上就是
\begin{align}
e_\ell &\leftarrow S_\ell^{\nu_1}(0,r_\ell),\\
\tilde r_\ell &\leftarrow r_\ell-A_\ell e_\ell,\\
r_{\ell+1} &\leftarrow R_\ell\tilde r_\ell.
\end{align}

到底以后调用哪个 OpenMP loop、GPU kernel、EB operator 或 ghost fill，不由 \texttt{mgVcycle()} 决定；它只调用 \texttt{linop} 提供的数值动作。

上升阶段同样直接：

\begin{verbatim}
addInterpCorrection(amrlev, mglev);
linop.smooth(amrlev, mglev,
             cor[amrlev][mglev],
             res[amrlev][mglev], false, nu2);
\end{verbatim}

对应
\begin{align}
e_\ell &\leftarrow e_\ell + P_\ell e_{\ell+1},\\
e_\ell &\leftarrow S_\ell^{\nu_2}(e_\ell,r_\ell).
\end{align}

因此 AMReX 的软件结构可以压缩为：
\begin{equation}
\boxed{
\texttt{MLMGT}
\text{ 保持 V-cycle 语义，}
\quad
\texttt{MLLinOpT}
\text{ 提供具体算子实现。}
}
\label{eq:software-amrex-driver-operator}
\end{equation}

这正是\chapref{ch:portability}所说的 numerical semantics 与 backend implementation 分离的真实工程例子。

\subsection{Bottom Solver 策略}

在最粗 MG level，\texttt{bottomSolve()} 最终进入
\texttt{actualBottomSolve()}。源码首先切换到

\begin{verbatim}
ParallelContext::push(linop.BottomCommunicator());
\end{verbatim}

然后根据策略分支：

\begin{itemize}
\item \texttt{BottomSolver::smoother}：继续调用 \texttt{linop.smooth}；
\item \texttt{BottomSolver::hypre}：调用 \texttt{bottomSolveWithHypre}；
\item \texttt{BottomSolver::petsc}：调用 \texttt{bottomSolveWithPETSc}；
\item \texttt{custom}：调用 operator 自己提供的 coarse solver；
\item 其他内部路径：使用 AMReX 的 CG/BiCGStab bottom solve。
\end{itemize}

这里有一个重要实现结论：

\begin{keybox}
在 AMReX MLMG 中选择 hypre/PETSc 作为 bottom solver，并不意味着整个 multigrid hierarchy 改成 hypre/PETSc。AMReX 仍然负责 fine/intermediate GMG levels，只在真正 bottom level 把线性系统交给外部求解器。
\end{keybox}

这类“外部库只接管底层”的组合模式，在大型 PDE 软件中比“整套 solver 全部替换”常见得多。

\section{AMReX 的粗层网格与所有权}

前几章把 agglomeration 和 consolidation 分别解释为
\textbf{任务粒度重组}与\textbf{所有权收缩}。在
\texttt{MLLinOpT::defineGrids()} 中，这两个概念可以直接从源码看到。

\subsection{Agglomeration 与 \texttt{BoxArray}}

AMReX 首先检查 coarse AMR level 是否适合 agglomeration，然后连续测试 MG coarsening。

其触发判据之一可以读成：
\begin{equation}
\frac{\text{coarsened domain points}}
{\text{number of boxes}}
<
\texttt{agg\_grid\_size}^{d}.
\end{equation}

触发以前，源码主要做
\begin{verbatim}
m_grids[0].push_back(coarsen(a_grids[0], ...));
m_dmap [0].push_back(a_dmap[0]);
\end{verbatim}

也就是继续沿用原 patch 拓扑和原 ownership。

触发以后，源码改为从更大的 coarse bounding box 重新生成
\texttt{BoxArray}，并用 \texttt{maxSize(agg\_grid\_size)} 控制 box 尺寸。

所以 agglomeration 的准确解释是：
\begin{equation}
\boxed{
\text{重写 coarse-level spatial task decomposition}.
}
\end{equation}

它首先解决的是 box 太碎、每 box work 太小的问题。

\subsection{Consolidation 与 \texttt{DistributionMapping}}

另一条路径先计算
\begin{equation}
\texttt{avg\_npts}
=
\frac{\text{fine-grid points}}
{\text{active ranks}}.
\end{equation}

随着 coarsening，如果估计的每 rank work 低于
\texttt{con\_grid\_size}$^d$ 阈值，就把相应 MG level 的
\texttt{DistributionMapping} 留空，后续再调用

\begin{verbatim}
makeConsolidatedDMap(...);
\end{verbatim}

生成新的 ownership mapping。

最后还有：

\begin{verbatim}
m_bottom_comm =
    makeSubCommunicator(m_dmap[0].back());
\end{verbatim}

这和\chapref{ch:mpi}的抽象完全一致：

\begin{equation}
\boxed{
\texttt{BoxArray}
=
\text{几何任务分块},
\qquad
\texttt{DistributionMapping}
=
\text{rank ownership}.
}
\label{eq:software-amrex-ba-dm}
\end{equation}

Agglomeration 主要改变前者，consolidation 主要改变后者。源码层面的这个区分，比“两个参数都能优化粗网格”更重要。

\section{hypre BoomerAMG 的层级构造}
\label{sec:software-hypre-source}

BoomerAMG 的源码组织与 AMReX 很不同。它没有已知几何 hierarchy，而是围绕
\texttt{hypre\_ParCSRMatrix} 数组构造整个代数层级。

最值得追的入口是
\begin{verbatim}
hypre_BoomerAMGSetup(...)
\end{verbatim}
所在的 \texttt{par\_amg\_setup.c}。

\subsection{强连接矩阵构造}

标准路径直接调用：

\begin{verbatim}
hypre_BoomerAMGCreateS(
    A_array[level],
    strong_threshold,
    max_row_sum,
    num_functions,
    dof_func_data,
    &S);
\end{verbatim}

这里的 \texttt{S} 就是式~\eqref{eq:software-amg-setup-chain}中的 strength graph。

因此 \texttt{strong\_threshold} 不是“某个神秘调参项”；它直接决定哪些矩阵边进入后续 coarsening graph。

\subsection{粗化策略}

源码随后根据 \texttt{coarsen\_type} 选择：

\begin{verbatim}
hypre_BoomerAMGCoarsenFalgout(...)
hypre_BoomerAMGCoarsenPMIS(...)
hypre_BoomerAMGCoarsenHMIS(...)
hypre_BoomerAMGCoarsenRuge(...)
...
\end{verbatim}

例如 \texttt{coarsen\_type == 8} 进入 PMIS，
\texttt{10} 进入 HMIS。

这揭示了 BoomerAMG 的一个重要架构事实：

\begin{equation}
\boxed{
\texttt{BoomerAMG}
\text{ 不是一个固定 AMG 算法，}
\quad
\text{而是一套可组合的 AMG setup policy family}.
}
\end{equation}

输出通常通过 \texttt{CF\_marker} 表示 coarse/fine splitting，后续 interpolation builder 会消费这个状态。

\subsection{插值构造}

在同一个 setup 文件中，
\texttt{interp\_type} 又会选择不同 builder，例如：

\begin{itemize}
\item \texttt{hypre\_BoomerAMGBuildMultipass}
\item \texttt{hypre\_BoomerAMGBuildInterpLS}
\item \texttt{hypre\_BoomerAMGBuildDirInterp}
\item \texttt{hypre\_BoomerAMGBuildExtPIInterp}
\item \texttt{hypre\_BoomerAMGBuildStdInterp}
\item \texttt{hypre\_BoomerAMGBuildInterpOnePnt}
\end{itemize}

最终都产生一个 \texttt{ParCSRMatrix *P}。

因此从源码看，AMG setup 的中间状态非常清楚：
\begin{equation}
A_\ell
\rightarrow
S_\ell
\rightarrow
\texttt{CF\_marker}_\ell
\rightarrow
P_\ell.
\label{eq:software-hypre-setup-state}
\end{equation}

\subsection{粗网格算子构造}

标准 Galerkin 路径最终调用

\begin{verbatim}
hypre_BoomerAMGBuildCoarseOperatorKT(
    P_array[level],
    A_array[level],
    P_array[level],
    keepTranspose,
    &A_H);
\end{verbatim}

或者走模块化的
\texttt{hypre\_ParCSRMatrixRAPKT}。

这不是教材里抽象的一个矩阵公式而已。它意味着 setup 真实地执行 distributed sparse matrix--matrix multiplication，并为新矩阵建立下一层通信结构。

所以 BoomerAMG 的 setup 成本来源至少包括：

\begin{equation}
\boxed{
\text{strength construction}
+
\text{parallel graph coarsening}
+
\text{interpolation construction}
+
\text{ParCSR RAP/SpGEMM}.
}
\label{eq:software-hypre-setup-cost}
\end{equation}

这也是为什么“solve 阶段 SpMV 很快”并不能推出动态重建 AMG hierarchy 也很便宜。

\section{hypre BoomerAMG 的 V-Cycle 状态机}

真正执行 cycle 的入口是
\begin{verbatim}
hypre_BoomerAMGCycle(amg_data, F_array, U_array)
\end{verbatim}
位于 \texttt{par\_cycle.c}。

它持有：

\begin{itemize}
\item \texttt{A\_array[level]}
\item \texttt{P\_array[level]}
\item \texttt{R\_array[level]}
\item \texttt{F\_array[level]}
\item \texttt{U\_array[level]}
\item 当前 \texttt{level}
\item \texttt{lev\_counter}
\item \texttt{cycle\_param}
\end{itemize}

和 PETSc 的递归实现不同，BoomerAMG 用这些控制变量显式决定“下一步去粗层、留在本层还是回细层”。

\subsection{下降阶段的 Residual 与限制}

向粗层移动时，源码先计算 residual：

\begin{verbatim}
hypre_ParCSRMatrixMatvecOutOfPlace(
    -1.0, A_array[fine_grid],
    U_array[fine_grid],
    1.0, F_array[fine_grid],
    Vtemp);
\end{verbatim}

也就是
\begin{equation}
V
=
F_\ell-A_\ell U_\ell.
\end{equation}

随后 restriction 直接由分布式稀疏矩阵乘实现。若 restriction 以 transpose 关系表达，会调用

\begin{verbatim}
hypre_ParCSRMatrixMatvecT(
    R_array[fine_grid],
    Vtemp,
    F_array[coarse_grid]);
\end{verbatim}

否则调用普通 \texttt{ParCSRMatrixMatvec}。

所以在 BoomerAMG 的 solve 阶段，
\textbf{restriction 不是某个特殊“AMG kernel”}，它本质上已经变成一个 ParCSR operator apply。

\subsection{上升阶段的延拓与校正}

从 coarse level 返回时：

\begin{verbatim}
hypre_ParCSRMatrixMatvec(
    1.0, P_array[fine_grid],
    U_array[coarse_grid],
    1.0, U_array[fine_grid]);
\end{verbatim}

对应
\begin{equation}
U_\ell
\leftarrow
U_\ell+P_\ell U_{\ell+1}.
\end{equation}

因此 BoomerAMG 的 solve 热路径可以概括成：

\begin{equation}
\boxed{
\text{smoother}
+
\text{ParCSR SpMV residual}
+
\text{ParCSR restriction}
+
\text{ParCSR prolongation}
}
\label{eq:software-hypre-solve-core}
\end{equation}

这也直接解释了 GPU 版本为什么非常依赖 device-native ParCSR matvec、relaxation 与 communication package。

\subsection{平滑器策略}

\texttt{par\_cycle.c} 中 smoother 会根据配置跳到不同路径，包括
Jacobi/\(\ell_1\)-Jacobi、Chebyshev、CG 类 smoother、FCF Jacobi、Schwarz 等。

从软件架构上看：

\begin{equation}
\boxed{
\text{cycle state machine 很稳定，}
\quad
\text{smoother policy 是最大的可变分支之一。}
}
\end{equation}

因此 profiler 若只把全部时间归到
\texttt{hypre\_BoomerAMGCycle}，信息仍然太粗；至少要继续拆到
Relaxation、Residual、Restriction、Interpolation 与 coarse solve。

\section{PETSc PCMG 的层级对象}
\label{sec:software-petsc-source}

PETSc 的实现思路与前两者再次不同。

\texttt{PC\_MG\_Levels} 直接把一个 level 的状态收在一个结构里。源码中的关键成员包括：

\begin{verbatim}
Vec b, x, r;
Mat A;
KSP smoothd;
KSP smoothu;
Mat interpolate;
Mat restrct;
\end{verbatim}

这里
\texttt{smoothd} 是下降阶段 smoother，
\texttt{smoothu} 是上升阶段 smoother。

因此 PETSc 没有把“smoother”硬编码成一个内部 relaxation enum，而是直接嵌入另一个完整的 \texttt{KSP} 对象。

这意味着一个 level 的抽象实际是：
\begin{equation}
\boxed{
\text{level}
=
(A,b,x,r)
+
(P,R)
+
(\text{pre-KSP},\text{post-KSP}).
}
\label{eq:software-petsc-level-object}
\end{equation}

\subsection{\texttt{PCMGMCycle\_Private()} 的递归 V-Cycle}

PETSc 的 multiplicative cycle 源码几乎可以当作标准递归 V-cycle 的参考实现。

先做 pre-smooth：
\begin{verbatim}
KSPSolve(mglevels->smoothd,
         mglevels->b,
         mglevels->x);
\end{verbatim}

再由函数指针计算 residual：
\begin{verbatim}
(*mglevels->residual)(
    mglevels->A,
    mglevels->b,
    mglevels->x,
    mglevels->r);
\end{verbatim}

再 restriction：
\begin{verbatim}
MatRestrict(mglevels->restrct,
            mglevels->r,
            mgc->b);
\end{verbatim}

随后直接递归：

\begin{verbatim}
while (cycles--)
    PCMGMCycle_Private(..., mglevelsin - 1, ...);
\end{verbatim}

返回以后：

\begin{verbatim}
MatInterpolateAdd(
    mglevels->interpolate,
    mgc->x,
    mglevels->x,
    mglevels->x);

KSPSolve(mglevels->smoothu,
         mglevels->b,
         mglevels->x);
\end{verbatim}

对照\chapref{ch:mg-cycles}可见，PETSc 基本没有隐藏 V-cycle 数学结构。

\subsection{PETSc 的对象组合架构}

从上面的源码可以看到：

\begin{itemize}
\item smooth 是 \texttt{KSP}；
\item restriction/prolongation 是 \texttt{Mat}；
\item state 是 \texttt{Vec}；
\item operator 是 \texttt{Mat}；
\item cycle 是 \texttt{PC} 内部控制流。
\end{itemize}

因此 PETSc 的性能可移植性很大程度上取决于这些底层对象的具体类型。例如同一个 \texttt{PCMGMCycle\_Private} 不需要知道矩阵最终由 CPU AIJ、CUDA、HIP 或其他 backend 执行。

这就是对象组合式架构的价值：
\begin{equation}
\boxed{
\text{cycle code 不直接依赖硬件，}
\quad
\text{backend 能力由 Mat/Vec/KSP 实现承担。}
}
\end{equation}

代价也很明确：若某个 backend 对 \texttt{MatRestrict}、某类 smoother 或 coarse solver 的 device path 不完整，统一的上层 cycle 本身并不能自动补救性能。

\section{PETSc PCGAMG 的层级构造}

PETSc 中另一个很值得源码学习的设计是：
\textbf{PCGAMG 并没有重新写一套 V-cycle executor。}

\texttt{PCSetUp\_GAMG()} 的 setup 循环会依次做：

\begin{verbatim}
PCGAMGCreateGraph(pc, Aarr[level], &Gmat);
pc_gamg->ops->coarsen(...);
pc_gamg->ops->prolongator(..., &Prol11);

if (pc_gamg->ops->optprolongator)
    pc_gamg->ops->optprolongator(...);

pc_gamg->ops->createlevel(
    ..., &Prol11, &Aarr[level1], ...);
\end{verbatim}

然后 hierarchy 构造完以后调用：

\begin{verbatim}
PCSetUp_MG(pc);
\end{verbatim}

这给出一个非常清晰的软件分层：

\begin{equation}
\boxed{
\texttt{PCGAMG}
=
\text{AMG hierarchy factory},
\qquad
\texttt{PCMG}
=
\text{multigrid execution engine}.
}
\label{eq:software-petsc-gamg-mg}
\end{equation}

这一点比“PETSc 同时支持 PCMG 和 GAMG”更值得记忆。

\subsection{GAMG 的粗空间策略}

\texttt{PCGAMGCreateGraph()} 本身只是：

\begin{verbatim}
pc_gamg->ops->creategraph(pc, A, G);
\end{verbatim}

具体实现由 \texttt{agg.c}、\texttt{classical.c}、
\texttt{geo.c} 等安装到函数表中。

同样地，
\texttt{coarsen}、
\texttt{prolongator}、
\texttt{optprolongator}、
\texttt{createlevel}
也是策略接口。

因此 PETSc GAMG 的“多种 AMG”不是在主循环里写大量硬编码分支，而是通过 operation table 插入不同 hierarchy construction policy。

\subsection{粗层进程收缩}

GAMG 的 coarse level 构造代码还会估计新的 active process 数：
\begin{equation}
P_{\mathrm{new}}
\approx
\frac{N_{\mathrm{coarse}}}
{\texttt{min\_eq\_proc}}.
\end{equation}

若 active process 数不变，路径可以直接执行
\begin{verbatim}
MatPtAP(Amat_fine, Pold, ..., A_coarse);
\end{verbatim}

若需要减少进程，则进入 repartition/reduce 路径，重新构造 coarse matrix 与 prolongator ownership。

这与 AMReX consolidation 在概念上类似，但数据对象完全不同：

\begin{itemize}
\item AMReX 重写 \texttt{DistributionMapping}；
\item PETSc GAMG 重写 coarse \texttt{Mat} / prolongator 的分布布局。
\end{itemize}

两者都在实现\chapref{ch:mpi}的同一个原则：
\begin{equation}
\boxed{
N_\ell/P_\ell
\text{ 太小时，coarse ownership 必须改变。}
}
\end{equation}

\section{AMReX、hypre 与 PETSc 的实现比较}

现在可以不依赖软件宣传页，直接比较三种设计。

{\small
\begin{longtable}{@{}>{\raggedright\arraybackslash}p{2.3cm}
                        >{\raggedright\arraybackslash}p{3.5cm}
                        >{\raggedright\arraybackslash}p{3.5cm}
                        >{\raggedright\arraybackslash}p{3.5cm}@{}}
\toprule
问题 & AMReX MLMG & hypre BoomerAMG & PETSc PCMG/GAMG\\
\midrule
\endfirsthead
\toprule
问题 & AMReX MLMG & hypre BoomerAMG & PETSc PCMG/GAMG\\
\midrule
\endhead
Hierarchy 来源 &
几何、BoxArray、coarsening &
ParCSR strength/coarsening/interpolation &
PCMG 可由 DM/用户提供；GAMG 从 Mat 构造\\

Level state &
\texttt{res/cor/sol} + operator hierarchy &
\texttt{A/P/R/U/F} arrays &
\texttt{PC\_MG\_Levels} 中的 Vec/Mat/KSP\\

Cycle owner &
\texttt{MLMGT::mgVcycle} &
\texttt{hypre\_BoomerAMGCycle} 状态机 &
\texttt{PCMGMCycle\_Private} 递归\\

Smoother &
\texttt{linop.smooth()} &
relax/smoother type dispatch &
每层独立 \texttt{KSP}\\

Restriction &
operator 专用实现 &
ParCSR matvec / matvecT &
\texttt{MatRestrict}\\

Prolongation &
operator 专用 interpolation &
ParCSR matvec-add &
\texttt{MatInterpolateAdd}\\

Coarse operator &
几何/离散算子生成 &
显式 ParCSR RAP &
PCMG 可外部提供；GAMG \texttt{MatPtAP}\\

粗层进程收缩 &
\texttt{DistributionMapping} consolidation + subcommunicator &
由 AMG hierarchy/distributed ParCSR 策略控制 &
GAMG createlevel 中 process reduction/repartition\\

底层求解 &
smoother、内部 CG/BiCGStab、hypre、PETSc、custom &
AMG coarse-level relaxation/solver path &
level 0 的 KSP/PC 配置\\
\bottomrule
\end{longtable}
}

三个项目的数学对象高度相似，但软件边界完全不同。

AMReX 的主线是
\begin{equation}
\boxed{
\text{geometry/operator polymorphism}
+
\text{explicit MG driver}.
}
\end{equation}

hypre BoomerAMG 的主线是
\begin{equation}
\boxed{
\text{ParCSR hierarchy construction}
+
\text{explicit level state machine}.
}
\end{equation}

PETSc 的主线是
\begin{equation}
\boxed{
\text{generic Mat/Vec/KSP composition}
+
\text{MG orchestration}.
}
\end{equation}

\section{源码级性能定位}

有了源码结构以后，性能问题可以落到具体阶段，而不是停留在“某库快六倍”。

\subsection{AMReX MLMG}

应优先 profile：

\begin{itemize}
\item \texttt{MLMG::mgVcycle\_down::<level>}
\item \texttt{MLMG::mgVcycle\_bottom}
\item smooth kernel
\item residual/operator apply
\item restriction/interpolation
\item agglomeration/consolidation 后的 coarse box/rank 数
\item bottom solver 构建与求解是否被重复执行
\end{itemize}

尤其要区分：
\begin{equation}
\text{MLMG cycle time}
\qquad\text{与}\qquad
\text{external bottom solver setup time}.
\end{equation}

\subsection{hypre BoomerAMG}

setup 应至少拆成：

\begin{equation}
\text{strength}
+
\text{coarsening}
+
\text{interpolation}
+
\text{RAP}
+
\text{communication-package setup}.
\end{equation}

solve 则至少拆为：

\begin{equation}
\text{relaxation}
+
\text{residual SpMV}
+
\text{restriction}
+
\text{interpolation}
+
\text{coarse solve}.
\end{equation}

若动态问题每次都重建 AMG hierarchy，单独比较 cycle time 会严重误导。

\subsection{PETSc PCMG/GAMG}

PETSc 最关键的是继续向对象内部追：

\begin{itemize}
\item 每层 \texttt{smoothd/smoothu} 到底是什么 KSP/PC；
\item \texttt{Mat} 类型是否驻留在期望 backend；
\item \texttt{MatPtAP} 是否占 setup 主导；
\item GAMG process reduction 在哪一层触发；
\item coarse KSP 是否使用了完全不同的 solver/backend；
\item \texttt{-log\_view} 中 MG level events 与 MPI collectives 的对应关系。
\end{itemize}

如果只看到 \texttt{PCApply} 很慢，还没有定位到真正瓶颈。

\section{MueLu、Ginkgo 与 AmgX 的源码入口}

本章不再用相同篇幅重复介绍另外三个库。真正继续源码阅读时，可以沿同一套问题进入。

\paragraph{MueLu.}
优先寻找 hierarchy、level 与 factory 之间的关系：谁生成 coarse map、aggregate、tentative/smoothed prolongator 与 coarse operator；再追 Tpetra/Kokkos 对象如何决定 execution/memory backend。它适合观察\textbf{factory-driven hierarchy construction}。

\paragraph{Ginkgo.}
优先寻找 Executor、LinOp、Multigrid level factory 与 distributed matrix/vector 的边界。它适合观察\textbf{operator object + executor} 如何把同一算法放到 Reference、OpenMP、CUDA、HIP、DPC++ 等后端，并进一步研究 mixed precision。

\paragraph{AmgX.}
优先从 solver configuration 到具体 AMG level/smoother/aggregation 实现，再追 CUDA kernel 与 distributed halo。它适合观察\textbf{GPU-first AMG} 如何把 setup 与 solve 都围绕设备执行组织。

真正值得比较的不是类名数量，而是它们分别把
\[
A_\ell,\quad P_\ell,\quad R_\ell,\quad
S_\ell,\quad x_\ell,\quad b_\ell
\]
放在哪里，由谁构造，由谁消费，以及 coarse ownership 怎样变化。

\section{多重网格源码阅读方法}

面对任何新的 multigrid 软件，建议固定按下面顺序追：

\begin{enumerate}
\item 找 public \texttt{solve/apply/setup} 入口；
\item 找 level state 的实际数据结构；
\item 找 cycle driver，确认递归还是显式状态机；
\item 找 pre/post smoother 的真实调用；
\item 找 residual/operator apply；
\item 找 restriction/prolongation 的真实 representation；
\item 找 coarse operator construction；
\item 找 bottom solver；
\item 找 coarse ownership / communicator / repartition 逻辑；
\item 最后才看配置参数和 backend kernel。
\end{enumerate}

这一顺序能防止源码阅读退化为“搜到很多函数名但不知道谁控制谁”。

\section{本章小结}

成熟多重网格软件之间最大的差异，不是是否都实现了 V-cycle，而是\textbf{它们把 V-cycle 的稳定语义放在哪一层，把可替换策略放在哪一层}。

AMReX 用 \texttt{MLMGT} 驱动 cycle，用 \texttt{MLLinOpT} 隔离几何与离散算子；它的 agglomeration/consolidation 直接分别落到 BoxArray 与 DistributionMapping。BoomerAMG 的 setup 是明确的
\[
A\rightarrow S\rightarrow\text{coarsening}\rightarrow P\rightarrow RAP
\]
流水线，solve 则是围绕 ParCSR arrays 的显式 level 状态机。PETSc PCMG 把每层表示为 Mat/Vec/KSP 的组合，而 PCGAMG 只负责制造 hierarchy，最终仍把执行交还 PCMG。

因此阅读求解器源码时，不应先问“这个库支持哪些参数”，而应先回答：
\[
\boxed{
\text{谁拥有 hierarchy？谁拥有 cycle？谁拥有 level state？}
}
\]
这三个问题一旦回答，性能、并行和可扩展性问题才有明确落点。

\section*{习题}

\subsection*{AMReX 源码追踪}
\begin{enumerate}[label=\arabic*.]
\item 从 \texttt{MLMGT::solve} 开始，画出
\texttt{solve -> oneIter -> miniCycle -> mgVcycle -> bottomSolve}
调用图，并标出 convergence check 位于哪一层。
\item 在 \texttt{mgVcycle} 中把每条关键调用分别映射为
smooth、residual、restriction、coarse solve、prolongation、post-smooth。
\item 阅读 \texttt{MLLinOpT::defineGrids}，指出 agglomeration 首次改变 \texttt{BoxArray} 的位置，以及 consolidation 首次改变 \texttt{DistributionMapping} 的位置。
\item 解释 \texttt{BottomCommunicator()} 为什么属于 ownership 层，而不是 stencil 数学层。
\item 对一个 profiler trace，把 \texttt{mgVcycle\_down}、bottom、up-cycle 与具体 operator kernels 对齐。
\end{enumerate}

\subsection*{hypre 源码追踪}
\begin{enumerate}[label=\arabic*.]
\item 在 \texttt{par\_amg\_setup.c} 中沿
\texttt{CreateS -> Coarsen -> BuildInterp -> BuildCoarseOperator}
追一遍默认配置，记录每一步产生的新数据结构。
\item 改变 \texttt{coarsen\_type}，找出 PMIS/HMIS/Falgout 分支的源码入口，说明哪些状态保持不变。
\item 在 \texttt{par\_cycle.c} 中解释
\texttt{level}、\texttt{lev\_counter}、\texttt{cycle\_param}
怎样替代递归调用栈。
\item 找出 residual、restriction 与 interpolation 的三个 ParCSR matvec 调用，分别写成数学公式。
\item 设计 profiler 标签，把 AMG setup 与 cycle solve 完全分离，避免把 RAP 时间算进迭代时间。
\end{enumerate}

\subsection*{PETSc 源码追踪}
\begin{enumerate}[label=\arabic*.]
\item 根据 \texttt{PC\_MG\_Levels} 画出一个 level object，说明
\texttt{A,b,x,r,smoothd,smoothu,interpolate,restrct}
之间的关系。
\item 沿 \texttt{PCMGMCycle\_Private} 手工画出递归调用栈，比较 V-cycle 与 W-cycle 时 \texttt{cycles} 的变化。
\item 在 \texttt{PCSetUp\_GAMG} 中追
\texttt{creategraph -> coarsen -> prolongator -> createlevel -> PCSetUp\_MG}
调用链，解释为什么 GAMG 可以复用 PCMG executor。
\item 找出 \texttt{MatPtAP} 与 process reduction/repartition 相邻的源码，说明 coarse operator construction 与 ownership change 为什么耦合。
\item 使用 PETSc profiling 输出，把 \texttt{PCSetUp\_GAMG}、\texttt{PCApply\_MG}、level smoother 与 MPI collective 分开计时。
\end{enumerate}

\subsection*{跨软件架构比较}
\begin{enumerate}[label=\arabic*.]
\item 对同一个三层 Poisson 问题，分别用 AMReX、hypre、PETSc 的内部对象描述
$A_\ell,b_\ell,x_\ell,r_\ell,P_\ell,R_\ell$。
\item 比较“AMReX 调用 hypre 作为 bottom solver”和“直接使用 BoomerAMG 作为整个 preconditioner”在 hierarchy ownership 上的根本区别。
\item 设计一个重复求解 benchmark，同时记录 setup、apply、bottom solve 与 coarse redistribution；说明三套软件中哪些阶段最可能被复用。
\item 选择 MueLu、Ginkgo 或 AmgX 之一，按本章十步源码阅读法补充一页实现笔记，不允许只引用用户文档。
\end{enumerate}
